\section{Introduction}
\label{sec:introduction}

Let \(X=(X_1,\ldots,X_n)\) be a sequence of independent uniform binary
random variables, and let \(Y\) be obtained by passing every coordinate
through a binary symmetric channel (BSC) with crossover probability
\(p\in[0,1/2]\). A Boolean summary \(f(X)\) retains one bit from the
input. Courtade and Kumar conjectured that the mutual information between
this summary and the complete noisy observation is maximized by retaining
one of the original coordinates~\cite{CK}. Equivalently,
\begin{equation}
 I(f(X);Y)\leq 1-h_2(p),
 \label{eq:intro-ck}
\end{equation}
where \(h_2\) is the binary entropy function. Equality is attained by any
\textit{dictator} defined as \(f(X)=X_i\) for some \(i \in \{1,2,\dots,n\}\).

The problem is different from maximizing correlation or noise
stability. The observation \(Y\) is an \(n\)-dimensional random vector,
and information carried jointly by its coordinates need not decompose
into coordinatewise contributions. For example, the parity of two noisy
bits can contain information that is absent from either noisy bit alone.
Consequently, results for two Boolean outputs~\cite{AGKN,PPM} or for sums
of coordinatewise mutual informations~\cite{JW} do not directly imply
\eqref{eq:intro-ck}. The output \(f(X)\) may also be biased, which removes
several symmetries available in the balanced case.

This paper establishes \eqref{eq:intro-ck} in full generality. The proof
uses the correlation parameter \(\rho=1-2p\). After monotone sorting, the
mutual information becomes an entropy functional of the posterior mean
\(T_\rho f\). Its derivative admits an edge decomposition. The central
step is to strengthen the entropy-constrained edge envelope by retaining
the squared angular displacement lost in a Cauchy--Schwarz comparison.
The pivotal sets of the same Boolean function supply the Fourier and
logarithmic-Sobolev control needed to pay for this correction. The final
argument is dynamical: wherever a function has more information than a
dictator, its advantage must increase with \(\rho\). Such an advantage
cannot return to its nonpositive value at \(\rho=1\).

We also study the natural vector-valued extension in which \(f(X)\)
contains \(k\) bits. The scalar theorem does not tensorize through the
chain rule because conditioning on preceding output bits destroys the
uniform input law. In fact, the coordinate benchmark is false for
sufficiently large fixed \(k\). We give explicit shortened-Hamming
counterexamples for \(k=8,9,10\), extend them to every \(k\geq8\), and
formulate a threshold conjecture for \(2\leq k\leq7\).

\subsection{Related Work}

Courtade and Kumar introduced the full-vector problem and proved several
structural reductions, including the monotone sorting operation used
below~\cite{CK}. The conjecture stimulated a sequence of Fourier,
functional-inequality, and information-geometric approaches. The
Friedgut--Kalai--Naor theorem identifies the rigidity of Boolean functions
whose Fourier mass is concentrated on the first two levels~\cite{FKN}.
Samorodnitsky established the conjecture in a dimension-free high-noise
regime and developed strengthened inequalities for the entropy of noisy
Boolean functions~\cite{Sam,SamEntropy}. Ordentlich, Shayevitz, and
Weinstein obtained sharper balanced high-noise bounds using Fourier
analysis and hypercontractivity~\cite{OSW}, while Yang and Wesel gave a
complementary calculus argument in a dimension-dependent neighborhood of
complete noise~\cite{YW}.

Yu developed a general \(\Phi\)-stability framework and subsequently
proved broad local-optimality and balanced finite-range results using
quantitative Fourier estimates~\cite{YuPhi,Yu,YuSharp}. Javanmard and
Woodruff proved the dictator bound for the sum of the individual-coordinate
mutual informations and sharpened the high-noise analysis~\cite{JW}; the
related Boolean-hull and moment relaxations are discussed in~\cite{Gemini}.
Li and M\'edard connected the problem to posterior moments and
noninteractive correlation distillation~\cite{LM}, and Barnes and
\"Ozg\"ur showed that the balanced conjecture is essentially equivalent
to a symmetrized Li--M\'edard formulation~\cite{BO}.

A second line of work seeks stronger entropy or isoperimetric statements.
Anantharam, Bogdanov, Chakrabarti, Jayram, and Nair proposed a Hellinger
inequality implying the conjecture~\cite{Hellinger}. Chen and Nair placed
this proposal in an ordered family of entropy inequalities and related
the low-noise limits to sharp cube isoperimetry~\cite{CN}. The critical
square-root isoperimetric inequality was recently established by Durcik,
Ivanisvili, Roos, and Xie~\cite{DIRX}. These limiting or stronger
statements do not themselves yield the finite-noise Shannon inequality.

The closest antecedent to the present proof is the differential-equation
framework odf Chen, Gohari, and Nair~\cite{CGN}. Their entropy-flow identity
and entropy-constrained two-point extremal formula are used here. 

Two
concurrent works establish the scalar conjecture in full
generality~\cite{CGJLMNW,KyTran}. Chen, Gohari, Javanmard, Lin, Mirrokni,
Nair, and Woodruff close the differential-equation program through an
unrestricted four-moment Bellman inequality for a hybrid potential,
using analytic reductions together with certified interval arithmetic.
Ky and Tran instead develop an entropy-production and spectral framework.
Their method is substantially closer to ours: both arguments use monotone
compression, edgewise entropy-production estimates, Fourier control of
functions with diffuse coordinate dependence, separate channel regimes,
and a differential contradiction at the noiseless endpoint. The
quantitative implementations differ. Our proof centers on an angular
edge correction and logarithmic-Sobolev estimates for pivotal sets,
whereas Ky and Tran use a profile-clock integration and selected-coordinate
spectral estimates. Finally, neither of~\cite{CGJLMNW,KyTran} considers the \(k\)-bit extension
studied in the present paper where we conjecture that up to $k=7,$ the most informative \(k\)-bit function follows a similar structure.

Related operator and channel perspectives include principal inertia
components~\cite{PIC} and mutual-information bounds based on adjacency
events~\cite{HOS}. Continuous analogues and stochastic-interpolation
variants were developed in~\cite{KOW,EMR}.

The present paper additionally develops the vector-valued extension in
which the quantizer reports \(k\) bits. For this problem, Chandar and
Tchamkerten characterized the positive-rate asymptotics and
gave Hamming- and Golay-code counterexamples with \(k\geq11\)~\cite{CT}.
Huleihel and Ordentlich proved coordinate optimality at the complementary
endpoint \(k=n-1\)~\cite{HO}. The finite-block behavior observed in our
code search is also related qualitatively to recent results showing that
structured Reed--Solomon-based codebooks can outperform random codebooks
under average Hamming covering radius~\cite{RMcover}; the objective in
that work is different from BSC mutual information.

\subsection{Contributions and Organization}

Our fist main contribution is Theorem~\ref{thm:CK}, which proves the
Courtade--Kumar conjecture in every dimension and for every crossover
probability, without any balance assumption on the Boolean function. The
proof follows a single dynamical strategy. For
$
 I_f(\rho)\coloneqq I(f(X);Y_\rho),
 \phi(\rho)\coloneqq I(X_1;Y_\rho),
 \Delta_f(\rho)\coloneqq I_f(\rho)-\phi(\rho),
$
we show that a positive advantage over a coordinate projection would have
to grow strictly as the channel becomes less noisy. Section~\ref{sec:mechanism}
introduces the posterior representation of \(I_f\), the required Fourier
notation, and the monotone reduction. In particular,
Lemma~\ref{lem:sorting} shows that coordinatewise sorting preserves any
advantage, allowing the remainder of the proof to exploit the
edge structure of a monotone Boolean decision set. An entropy-flow
identity expresses \(\rho I_f'(\rho)\) as an action functional of the
posterior mean \(T_\rho f\). Section~\ref{sec:edges} then decomposes this
action over edges of the discrete cube. The key estimate,
Lemma~\ref{lem:edges}, strengthens the usual two-point entropy envelope by
retaining a squared angular correction that is exact for dictators. The
rest of the proof shows that the Boolean and spectral structure of \(f\)
always provides enough reserve to pay for this correction.

Writing \(\ell=\operatorname{arctanh}\rho\), we divide the channel into
the small-correlation range \(0\leq\ell\leq31/20\), the intermediate
range \(31/20\leq\ell\leq7/2\), and the near-noiseless tail
\(\ell\geq7/2\). Equivalently, the two junctions are
\(\rho=\tanh(31/20)\) and \(\rho=\tanh(7/2)\). These cutoffs are technical
junctions at which the corresponding estimates meet, rather than
asserted phase transitions in the problem. Proposition~\ref{prop:low}
settles the small-correlation range directly, at every output bias, by
combining a Fourier energy gap with scalar entropy bounds. In the
intermediate range, the pivotal sets of \(f\) supply the needed reserve:
logarithmic Sobolev inequalities control the noisy entropies of their
indicators, and a log-sum argument combines the coordinatewise costs with
the global edge budget. This yields
Proposition~\ref{prop:middle-growth}, according to which
\(\Delta_f(\rho)>0\) implies \(\Delta_f'(\rho)>0\) throughout that range.
The bounded angular reserve is insufficient by itself near the noiseless
endpoint, so the tail argument supplements the same edgewise inequality
with degree-weighted Fourier information. Lemma~\ref{lem:deficit}
identifies the relevant coordinate deficits,
Lemma~\ref{lem:baseline} converts the information advantage into an
entropy-funded spectral baseline, and Lemma~\ref{lem:budget} turns the
resulting budget comparison into strict information growth.
Proposition~\ref{prop:local} separately rules out a counterexample having
a sufficiently influential coordinate. Section~\ref{sec:coverage}
assembles these complementary estimates in
Lemma~\ref{lem:coverage}, which gives the global implication
\[
 \Delta_f(\rho)>0\quad\Longrightarrow\quad\Delta_f'(\rho)>0.
\]
Since \(\Delta_f(1)=H(f(X))-\log 2\leq0\), a positive advantage cannot
persist to the noiseless endpoint; this contradiction completes the proof
of Theorem~\ref{thm:CK}.

Our second contribution, developed in
Section~\ref{sec:multibit-extension}, concerns the natural extension in
which the quantizer reports \(k\) bits and the coordinate benchmark is
\(k(1-h_2(p))\). Proposition~\ref{prop:affine-multibit} shows that affine
linear maps cannot violate this benchmark. Nonlinear nearest-codeword
quantizers built from linear codebooks behave differently: an exact
high-noise expansion, expressed through the coset-leader distribution,
produces shortened-Hamming counterexamples for \(k=8,9,10\) in
Proposition~\ref{prop:k8-k10-counterexamples}. A direct-sum construction
in Proposition~\ref{cor:failure-all-k-ge-8} extends the failure to every
\(k\geq8\). For \(k<8\), exhaustive searches over several structured
code families and Monte Carlo experiments with random codebooks find no
violation, while also showing a clear advantage of structured codebooks
over random ones in this finite-block regime. These findings motivate
Conjecture~\ref{conj:seven-bit-threshold}, which predicts that coordinate
projections remain optimal for \(2\leq k\leq7\). The appendices provide
the scalar inequalities, logarithmic-Sobolev and Fourier estimates,
channel-range comparisons, endpoint reductions, exact arithmetic
certificates, and reproducible numerical diagnostics used throughout the
proof.

